\documentclass[letterpaper]{article}

\usepackage[margin=1in]{geometry}
\usepackage{amsmath,amssymb}
\usepackage{booktabs}
\usepackage{multirow}
\usepackage{graphicx}
\usepackage{xcolor}
\usepackage{hyperref}
\usepackage{authblk}
\usepackage{caption}
\usepackage{subcaption}
\usepackage{enumitem}

\title{Why Token Pruning Crashes Some Robot Policies but Not Others:\\The Role of Vision Connectors in VLA Token Compressibility}

\author[1]{Anonymous Authors}
\affil[1]{Anonymous Institution}
\date{}

\begin{document}
\maketitle

%% ===== ABSTRACT =====
\begin{abstract}
Vision-Language-Action (VLA) models encode visual observations as hundreds of tokens, creating a computational bottleneck for real-time robot control. A growing body of work proposes inference-time token pruning to address this, yet these methods exhibit wildly inconsistent results across architectures---succeeding on some VLAs while catastrophically failing on others. We present the first systematic analysis of \emph{why} token pruning behavior differs across VLA architectures. Testing 5 VLA models (450M--7B parameters) spanning 3 vision connector families on the LIBERO manipulation benchmark, we uncover a clean taxonomy: pure linear projection ($\pi$0) is pruning-immune (loss improves 2.4\% at 90\% removal); linear projection with deep post-processing (X-VLA, 12-layer BART encoder) exhibits a threshold-cliff pattern (97.6\% SR at $k{=}0.9$, crashing to 56.2\% at $k{=}0.7$); multi-layer MLP connectors (OpenVLA) degrade smoothly ($-$49\% relative at 50\% removal); and pixel-shuffle connectors (SmolVLA) collapse catastrophically (0\% SR at 5\% removal). We show that the coefficient of variation (CV) of token importance scores---computable from a single forward pass---forms a continuous predictor of this spectrum (CV: $0.038 \to 0.151 \to 0.215 \to 0.314$). For the hardest case (pixel-shuffle), we propose Vision Token Re-encoding (VTR), a lightweight cross-attention module ($\sim$22M parameters) that recovers from 0\% to 60\%+ SR at 75\% token compression. VTR's effectiveness is itself connector-dependent: it specifically solves pixel-shuffle spatial entanglement, while degrading MLP-projector models (0\% SR) where tokens are already independently meaningful. This architecture-dependency of both the problem and the fix reinforces our core finding that the vision connector is the primary determinant of token compressibility in VLAs.
\end{abstract}

%% ===== 1. INTRODUCTION =====
\section{Introduction}

Vision-Language-Action (VLA) models have emerged as a promising paradigm for generalist robot manipulation, leveraging pre-trained vision-language backbones to map visual observations directly to actions~\cite{brohan2023rt2,kim2024openvla,black2024pi0}. However, these models encode each visual observation as hundreds of tokens (typically 128--256), which dominate the sequence length and create a computational bottleneck during autoregressive action decoding. This motivates a growing line of work on inference-time vision token reduction.

Recent methods---including attention-based pruning, token merging, and importance-weighted selection---have demonstrated promising results on specific VLA architectures. However, a critical question remains unanswered: \emph{why do these methods succeed on some architectures but catastrophically fail on others?} The implicit assumption that token pruning behavior is architecture-independent has never been systematically tested.

We address this gap through a comprehensive cross-architecture analysis of token pruning robustness. Testing 5 VLA models spanning 3 vision connector families (linear projection, multi-layer MLP, and pixel-shuffle) and 4 model sizes (450M--7B), we discover that the \emph{effective vision-to-LLM projection pipeline}---not model size, not the vision encoder, and not the action head---determines pruning behavior. Our findings reveal a clean taxonomy:

\begin{enumerate}[leftmargin=*,nosep]
    \item \textbf{Pure linear projection} ($\pi$0, 3.3B): pruning-immune. Removing 90\% of tokens \emph{improves} loss by 2.4\%.
    \item \textbf{Linear + deep post-processing} (X-VLA, 880M): threshold-graded. Robust at $k{\geq}0.9$ (97.6\% SR), cliff at $k{=}0.7$ (56.2\%), then graded degradation.
    \item \textbf{Multi-layer MLP} (OpenVLA, 7B): smooth graded degradation. 50\% removal yields $-$49\% relative SR.
    \item \textbf{Pixel-shuffle} (SmolVLA, 450M): catastrophic cliff-edge. Even 5\% removal causes 0\% task success.
\end{enumerate}

A key controlled comparison isolates the connector as causal: $\pi$0 (256 tokens, linear) and OpenVLA (256 tokens, 3-layer MLP) share the same token count and SigLIP encoder lineage. At $k{=}0.1$, $\pi$0 \emph{improves} 2.4\% while OpenVLA degrades 61\%---the only difference is the connector.

We further show that the coefficient of variation (CV) of token importance scores, computable in a single forward pass, forms a continuous predictor of degradation onset: CV $= 0.038$ (immune) $\to 0.151$ (threshold-graded) $\to 0.215$ (graded) $\to 0.314$ (catastrophic). For the most vulnerable architecture (pixel-shuffle), we propose Vision Token Re-encoding (VTR), which recovers from 0\% to 60\%+ SR at 75\% compression. Crucially, VTR's effectiveness is itself connector-dependent: it specifically repairs pixel-shuffle spatial entanglement but causes overfitting on MLP-projector models, reinforcing the centrality of the vision connector in VLA token compressibility.

%% ===== 2. BACKGROUND =====
\section{Background and Related Work}

\subsection{Vision-Language-Action Models}

VLA models adapt pre-trained vision-language models (VLMs) for robotic control by adding action prediction capabilities. The vision processing pipeline in a VLA typically consists of: (1) a frozen vision encoder (e.g., SigLIP~\cite{zhai2023sigmoid}) that produces patch embeddings, (2) a \emph{vision connector} that projects these embeddings into the LLM's token space, and (3) the LLM backbone that processes the resulting tokens alongside language instructions. The choice of vision connector varies substantially across architectures:

\begin{itemize}[leftmargin=*,nosep]
    \item \textbf{Linear projection}: $\pi$0~\cite{black2024pi0} uses a single linear layer, preserving the independence of individual patch embeddings.
    \item \textbf{Linear + deep post-processing}: X-VLA~\cite{zhen2024xvla} applies a linear projection followed by a 12-layer BART encoder, introducing substantial cross-token mixing.
    \item \textbf{Multi-layer MLP}: OpenVLA~\cite{kim2024openvla} uses a 3-layer MLP with GELU activations, creating partial dependencies through non-linear mixing.
    \item \textbf{Pixel-shuffle}: SmolVLA~\cite{smolvla2025} uses spatial folding to compress 1024 patches into 64 tokens per camera via channel-to-spatial reorganization, creating extreme cross-token entanglement.
\end{itemize}

\subsection{Token Efficiency in VLMs and VLAs}

Several concurrent works address vision token reduction for VLAs. Compressor-VLA, SpecPrune-VLA, LightVLA~\cite{lightvla2025}, and See What Matters~\cite{seewhatmatters2025} all propose pruning on MLP-projector VLAs. EfficientVLA reports 74.8\% $\to$ 74.2\% under 78\% pruning on CogACT (MLP projector)---consistent with our finding that MLP connectors tolerate pruning. However, none of these works examine \emph{why} pruning behavior differs across architectures or test on pixel-shuffle connectors. VLA-IAP's 97.8\% on OpenVLA is also consistent with MLP projector tolerance.

%% ===== 3. EXPERIMENTAL SETUP =====
\section{Experimental Setup}

\subsection{Models}

We evaluate 5 VLA models spanning 3 connector families and 4 model sizes: $\pi$0 (3.3B, linear), X-VLA (880M, linear + BART 12L), OpenVLA (7B, 3-layer MLP), OpenVLA-OFT (7B, 3-layer MLP, different action head), and SmolVLA (450M, pixel-shuffle). This coverage eliminates the model-size confound: $\pi$0 (3.3B, linear) is pruning-immune while OpenVLA (7B, MLP) degrades substantially---the larger model is \emph{more} vulnerable. X-VLA (880M) is also robust despite being 4$\times$ smaller than $\pi$0.

\subsection{Benchmark}

All closed-loop evaluations use LIBERO~\cite{liu2024libero}, a simulated manipulation benchmark with 10 tasks and 50 evaluation episodes per task (500 total). We report success rate (SR) as the primary metric.

\subsection{Pruning Protocol}

For each model, we evaluate at retention ratios $k \in \{0.95, 0.9, 0.7, 0.5, 0.3\}$, where $k{=}0.9$ means 90\% of tokens are retained. We use two pruning strategies: (1) \textbf{Random removal}: tokens are selected uniformly at random for removal (OpenVLA primary); (2) \textbf{L2-norm importance}: tokens are ranked by L2-norm and the lowest-norm tokens are removed (X-VLA primary, and OpenVLA comparison). This cross-method protocol strengthens findings: OpenVLA survives even random removal (the harder test), while X-VLA cliffs even with importance-based selection (the easier test).

\subsection{Diagnostic Metrics}

We compute two zero-cost diagnostic metrics from a single forward pass:

\textbf{Importance CV}: The coefficient of variation of L2-norm importance scores across vision tokens. Higher CV indicates more heterogeneous importance distribution and less tolerance to pruning.

\textbf{Attention entropy}: Mean entropy of suffix-to-vision attention weights across LLM layers. High entropy with low CV ($\pi$0: 0.959/0.038) indicates uniform, independent tokens. Low entropy with moderate CV (OpenVLA: 0.661/0.215) indicates selective attention with partial dependencies.

\subsection{Vision Token Re-encoding (VTR)}

VTR is a lightweight cross-attention module ($\sim$22M parameters) trained to re-encode vision tokens. $M$ learnable query tokens attend to the original $N$ vision tokens via cross-attention, producing a compressed representation. A decoder reconstructs $N$ tokens from $M$ queries. At $M{=}128$ (passthrough), no compression occurs---any performance gain isolates the re-encoding effect from compression. We sweep $M \in \{16, 32, 64, 96, 128\}$.

%% ===== 4. MAIN RESULTS =====
\section{Main Results}

\subsection{Four-Class Taxonomy of Pruning Robustness}

Table~\ref{tab:taxonomy} presents our central finding: a taxonomy of token pruning robustness determined by the vision connector architecture.

\begin{table}[t]
\centering
\caption{Cross-architecture token pruning taxonomy. Five VLA models spanning 3 connector families. Normalized loss (offline) and closed-loop LIBERO success rates (SR) at key retention ratios. CV = coefficient of variation of token importance.}
\label{tab:taxonomy}
\small
\begin{tabular}{@{}llccccc@{}}
\toprule
\textbf{Model} & \textbf{Connector} & \textbf{CV} & \textbf{$k{=}0.9$ SR} & \textbf{$k{=}0.7$ SR} & \textbf{$k{=}0.5$ SR} & \textbf{Class} \\
\midrule
$\pi$0 (3.3B) & Pure linear & 0.038 & --- & --- & 0.983$^*$ & Immune \\
X-VLA (880M) & Linear + BART 12L & 0.151 & 97.6\% & 56.2\% & 42.6\% & Threshold-graded \\
OpenVLA (7B) & 3-layer MLP & 0.215 & 62.0\% & 49.0\% & 37.4\% & Smooth graded \\
OFT (7B) & 3-layer MLP & --- & --- & --- & --- & (entropy $\approx$ OpenVLA) \\
SmolVLA (450M) & Pixel-shuffle & 0.314 & 0\% & 0\% & 0\% & Catastrophic \\
\bottomrule
\multicolumn{7}{l}{\footnotesize $^*$Normalized loss (offline metric); $k{=}0.1$ = 0.976 (improves 2.4\%).}
\end{tabular}
\end{table}

The taxonomy reveals a monotonic relationship between connector complexity, importance CV, and pruning vulnerability. The key controlled comparison---$\pi$0 (256 tokens, linear) vs.\ OpenVLA (256 tokens, 3-layer MLP)---isolates the connector as the causal variable: same token count, same encoder lineage, divergent pruning outcomes.

\subsection{Catastrophic: Pixel-Shuffle Cliff-Edge (SmolVLA)}

SmolVLA's pixel-shuffle connector spatially folds 1024 SigLIP patches into 64 tokens per camera (16$\times$ compression). Each resulting token encodes a 32$\times$32 pixel region via channel-to-spatial reorganization, creating extreme cross-token spatial entanglement.

The consequence is catastrophic: removing even 5\% of tokens ($k{=}0.95$, 7 of 128 tokens) yields 0\% success rate. Table~\ref{tab:smolvla_cliff} shows that all pruning strategies---random, attention-based (ACIS), token merging (ToMe), and geometric (Sobel edge)---produce identically zero performance across all retention ratios.

\begin{table}[t]
\centering
\caption{SmolVLA pruning results on LIBERO Object suite. Baseline SR = 44.8\%. All configurations yield 0\% SR regardless of pruning method or retention ratio.}
\label{tab:smolvla_cliff}
\small
\begin{tabular}{@{}lcccc@{}}
\toprule
\textbf{Method} & $k{=}0.9$ & $k{=}0.7$ & $k{=}0.5$ & $k{=}0.95$ \\
\midrule
Random       & 0\% & 0\% & 0\% & 0\% \\
ACIS         & 0\% & 0\% & 0\% & 0\% \\
ToMe         & 0\% & 0\% & 0\% & --- \\
Sobel (geometric) & 0\% & 0\% & 0\% & --- \\
\bottomrule
\end{tabular}
\end{table}

The official SmolVLA checkpoint exhibits the same cliff-edge behavior: Object suite (59.2\% baseline) drops to 0\% at $k{=}0.5$; Spatial suite (43.8\% baseline) drops to 0.6\% (3/500 episodes, 95\% CI [0.2\%, 1.7\%]).

\subsection{Smooth Graded: MLP Degradation (OpenVLA)}

OpenVLA's 3-layer MLP connector (with GELU activations) creates partial token dependencies through non-linear mixing, resulting in smooth, graded degradation as more tokens are removed. Table~\ref{tab:openvla_curve} presents the complete degradation curve.

\begin{table}[t]
\centering
\caption{OpenVLA pruning degradation curve on LIBERO (baseline SR = 74.0\%). Random pruning primary; L2-norm comparison shows near-identical curves ($\Delta < 3$pp), confirming pruning \emph{method} is not the determining factor.}
\label{tab:openvla_curve}
\small
\begin{tabular}{@{}cccccc@{}}
\toprule
$k$ & \textbf{Tokens Removed} & \textbf{Random SR} & \textbf{L2-norm SR} & \textbf{$\Delta$} & \textbf{Rel.\ Degradation} \\
\midrule
1.0  & 0\%  & 74.0\% & 74.0\% & 0    & --- \\
0.95 & 5\%  & 69.4\% & 70\%   & +0.6pp & $-$6\% \\
0.9  & 10\% & 62.0\% & 63\%   & +1pp   & $-$16\% \\
0.7  & 30\% & 49.0\% & 51\%   & +2pp   & $-$34\% \\
0.5  & 50\% & 37.4\% & 40\%   & +3pp   & $-$49\% \\
0.3  & 70\% & 20.6\% & 23\%   & +2pp   & $-$72\% \\
\bottomrule
\end{tabular}
\end{table}

The near-overlap of random and L2-norm pruning curves ($\Delta < 3$pp across all $k$) confirms that \emph{connector type, not pruning strategy, determines degradation behavior}. The maximum L2-norm advantage occurs at $k{=}0.5$ (+3pp), within the statistical noise of 500 evaluation episodes.

Seed replication confirms the smooth pattern's reproducibility: at $k{=}0.7$, seed42 = 49.0\% and seed43 = 47\% ($\Delta = 2$pp); at $k{=}0.5$, seed42 = 37.4\% and seed43 = 36\% ($\Delta = 1.4$pp). This low seed variance contrasts sharply with X-VLA's cliff-point variance (Section~\ref{sec:xvla}).

OpenVLA-OFT uses the same 3-layer MLP connector but a different action head. Its attention entropy (0.667) is virtually identical to standard OpenVLA (0.661), confirming that the connector---not the action head---determines the attention structure over vision tokens.

\subsection{Immune and Threshold-Graded: Linear Projection ($\pi$0 and X-VLA)}
\label{sec:xvla}

\textbf{Pure linear projection ($\pi$0).} $\pi$0's single linear layer preserves the independence of SigLIP patch embeddings. At $k{=}0.1$ (removing 90\% of tokens), the normalized loss is 0.976---a 2.4\% \emph{improvement}. Attention entropy is 0.959 (near-uniform across all 18 LLM layers, range [0.931, 0.983]), and importance CV = 0.038 (the lowest in our study). Every token carries nearly equal importance; removing any subset barely affects the LLM's processing.

\textbf{Linear + BART encoder (X-VLA).} X-VLA applies a linear projection followed by a 12-layer BART encoder before the LLM decoder. Despite the ``linear projection'' label, this deep post-processing introduces substantial cross-token mixing, elevating CV to 0.151 and producing a distinctive threshold-graded pattern. Table~\ref{tab:xvla_curve} presents the full degradation curve.

\begin{table}[t]
\centering
\caption{X-VLA L2-norm pruning curve on LIBERO (baseline SR = 99.2\%). Threshold at $k{=}0.7$: 97.6\% $\to$ 56.2\% ($-$41pp). Seed44 replication confirms cliff pattern with expected high variance at the threshold.}
\label{tab:xvla_curve}
\small
\begin{tabular}{@{}cccccc@{}}
\toprule
$k$ & \textbf{Seed42 SR} & \textbf{Seed44 SR} & \textbf{Mean} & \textbf{Pattern} \\
\midrule
1.0  & 99.2\% & $\sim$99\%  & $\sim$99\%  & Baseline \\
0.95 & 99.0\% & $\sim$98\%  & $\sim$98.5\% & Robust \\
0.9  & 97.6\% & 96.5\%     & $\sim$97\%  & Robust \\
0.7  & 56.2\% & 52\%       & $\sim$54\%  & \textbf{Cliff} ($-$43pp) \\
0.5  & 42.6\% & ---        & ---         & Below cliff \\
0.3  & 23.0\% & ---        & ---         & Graded degradation \\
0.2  & $\sim$4\% & ---     & ---         & Near-zero \\
0.1  & $\sim$0\% & ---     & ---         & Failure \\
\bottomrule
\end{tabular}
\end{table}

The cliff at $k{=}0.7$ reflects the BART encoder's partial token dependencies: removing 30\% of tokens crosses a threshold where critical inter-token information is irreversibly lost. Seed variance is characteristically high at the cliff point (seed42: 56.2\%, seed44: 52\%) but low in the robust region (seed42: 97.6\%, seed44: 96.5\% at $k{=}0.9$), consistent with a phase-transition-like phenomenon where small perturbations near the threshold produce large outcome differences.

\textbf{Key insight}: ``Linear projection'' is not sufficient for pruning immunity. The \emph{effective} pipeline matters: $\pi$0's pure linear projection yields CV = 0.038 and immunity, while X-VLA's linear projection + 12-layer BART encoder yields CV = 0.151 and a cliff at $k{=}0.7$.

\subsection{VTR Overcomes Pixel-Shuffle Cliff-Edge}
\label{sec:vtr}

For SmolVLA's catastrophic cliff-edge, we propose Vision Token Re-encoding (VTR): a lightweight cross-attention module that re-encodes the entangled pixel-shuffle tokens into a compressed, disentangled representation. Table~\ref{tab:vtr_msweep} presents the M-sweep results.

\begin{table}[t]
\centering
\caption{SmolVLA VTR M-sweep on LIBERO Object suite (3-seed, baseline = 44.8\%, hard pruning = 0\%). The M-curve is flat: all M values within $\pm$5pp, confirming re-encoding---not compression ratio---drives recovery.}
\label{tab:vtr_msweep}
\small
\begin{tabular}{@{}ccccccc@{}}
\toprule
$M$ & \textbf{Compression} & \textbf{Seed42} & \textbf{Seed43} & \textbf{Seed44} & \textbf{Mean$\pm$Std} \\
\midrule
16  & 87.5\% & 55.8\% & 58.4\% & 56.5\% & 56.9 $\pm$ 1.3\% \\
32  & 75\%   & 61.6\% & 59.2\% & 59.6\% & 60.1 $\pm$ 1.3\% \\
64  & 50\%   & 61.6\% & 60.0\% & 64.4\% & 62.0 $\pm$ 2.3\% \\
96  & 25\%   & 52.6\% & 58.2\% & 52.8\% & 54.5 $\pm$ 3.2\% \\
128 & 0\% (passthrough) & 59.8\% & 54.0\% & 60.2\% & 58.0 $\pm$ 3.5\% \\
\bottomrule
\end{tabular}
\end{table}

The M-curve is remarkably flat: M=16 (87.5\% compression, 56.9$\pm$1.3\%) $\approx$ M=128 (no compression, 58.0$\pm$3.5\%), with the difference $< 1.1$pp. This flatness demonstrates that the gain comes from \emph{re-encoding}---disentangling the spatially folded tokens---rather than any particular compression ratio. VTR achieves near-lossless compression at any ratio for pixel-shuffle connectors.

The M=128 passthrough control is critical: at M=128, VTR performs no compression but still re-encodes all tokens. The resulting 58.0\% SR (vs.\ 44.8\% baseline and 0\% hard pruning) isolates the re-encoding mechanism as the causal factor.

\textbf{Architecture independence of re-encoding.} An MLP adapter ($\sim$22.14M params, matching VTR's parameter count) achieves 63.4$\pm$2.6\% at M=32 vs.\ cross-attention VTR's 60.1$\pm$1.3\% (two-sample $t$-test: $p \approx 0.15$, not significant). Both architectures recover from 0\% to 60\%+, confirming that \emph{any} learned re-encoding suffices.

\textbf{Cross-suite generalization.} VTR trained on Object suite transfers: Spatial suite recovers from 0.6\% (pruning cliff) to 53.2\% at M=32; Goal suite recovers from 0\% to 56.4\%.

%% ===== 5. ANALYSIS =====
\section{Analysis}

\subsection{VTR is Connector-Dependent: MLP Projector Failure}

A central question is whether VTR provides universal compression or is specific to pixel-shuffle connectors. We test VTR on OpenVLA (3-layer MLP connector) across multiple configurations. Table~\ref{tab:vtr_openvla} summarizes the results.

\begin{table}[t]
\centering
\caption{VTR and VTC-only on OpenVLA (MLP connector, baseline = 74.0\%). VTR fails catastrophically on MLP-projector models---tokens are already independently meaningful, and re-encoding introduces destructive overfitting.}
\label{tab:vtr_openvla}
\small
\begin{tabular}{@{}llccc@{}}
\toprule
\textbf{Architecture} & $M$ & \textbf{SR} & \textbf{vs.\ Baseline} & \textbf{Failure Mode} \\
\midrule
VTR (full) & 128 & 0\% & $-$100\% & Identical action every step \\
VTC-only   & 32  & 7\% & $-$91\%  & Severe degradation \\
VTC-only   & 64  & 4\% & $-$95\%  & Severe degradation \\
VTC-only   & 96  & 2\% & $-$97\%  & Severe degradation \\
\bottomrule
\end{tabular}
\end{table}

VTR at M=128 on OpenVLA yields 0\% SR (0/44 episodes), with the model collapsing to output an identical action at every timestep---a signature of severe overfitting. The VTC-only variants (compressor without decoder) fare marginally better but remain catastrophically below baseline: M=32 achieves 7\%, M=64 achieves 4\%, and M=96 achieves 2\%.

The inverse M-curve on OpenVLA (larger M $\to$ worse performance) contrasts with SmolVLA's flat M-curve, and the explanation is mechanistic: MLP-projector tokens are already individually meaningful (CV = 0.215, moderate heterogeneity but no spatial entanglement). The cross-attention bottleneck forces 256 independent tokens through unnecessary re-encoding, and the additional $\sim$22M parameters overfit rapidly on LIBERO's small-scale demonstration data ($\sim$50 demos per task). For pixel-shuffle tokens, this same re-encoding disentangles spatial folding; for MLP tokens, it is a destructive operation.

This finding strengthens the paper's core narrative: even the \emph{fix} for token compressibility is architecture-dependent.

\subsection{Diagnostic Metrics: CV as Continuous Predictor}
\label{sec:diagnostic}

Table~\ref{tab:diagnostics} presents the diagnostic metrics across all models. Importance CV forms a continuous spectrum that predicts pruning behavior without requiring closed-loop evaluation.

\begin{table}[t]
\centering
\caption{Diagnostic metrics and pruning behavior across 5 VLA models. CV forms a continuous predictor: lower CV $\to$ later degradation onset. Entropy provides mechanistic context but is not independently predictive (SmolVLA: high entropy but catastrophic).}
\label{tab:diagnostics}
\small
\begin{tabular}{@{}llccl@{}}
\toprule
\textbf{Model} & \textbf{Connector} & \textbf{Entropy} & \textbf{CV} & \textbf{Behavior} \\
\midrule
$\pi$0 (3.3B)   & Pure linear        & 0.959 & 0.038 & Immune \\
X-VLA (880M)     & Linear + BART 12L  & 0.614 & 0.151 & Threshold-graded ($k{=}0.7$ cliff) \\
OpenVLA (7B)     & 3-layer MLP        & 0.661 & 0.215 & Smooth graded \\
OFT (7B)         & 3-layer MLP        & 0.667 & ---   & $\approx$ OpenVLA \\
SmolVLA (450M)   & Pixel-shuffle      & 0.952 & 0.314 & Catastrophic cliff \\
\bottomrule
\end{tabular}
\end{table}

CV captures the degree of non-linear mixing in the vision-to-LLM pathway: more mixing $\to$ higher CV $\to$ less tolerance to pruning. Notably, entropy alone is insufficient: SmolVLA has near-maximal entropy (0.952) similar to $\pi$0 (0.959), yet exhibits opposite pruning behavior. The distinction is that SmolVLA's uniform attention arises from spatial entanglement (every token mixes information from all spatial locations), while $\pi$0's uniform attention arises from true token independence.

X-VLA's CV = 0.151 correctly predicted its intermediate behavior before closed-loop evaluation. Its entropy (0.614) is close to OpenVLA's (0.661)---both go through non-linear mixing (BART encoder and MLP respectively)---providing mechanistic grounding for why both show degradation beyond a threshold, while differing in the sharpness of the cliff.

\subsection{Cross-Model Pruning Method Comparison}

To ensure that our taxonomy reflects connector properties rather than pruning method artifacts, we evaluate L2-norm importance pruning on OpenVLA alongside the primary random pruning results. Table~\ref{tab:openvla_curve} (Section 4.3) shows the complete comparison.

The two methods produce near-identical degradation curves: the maximum difference is 3pp at $k{=}0.5$, well within the statistical margin of 500-episode evaluation. At the key cross-model comparison point of $k{=}0.9$, OpenVLA L2-norm achieves 63\% vs.\ X-VLA L2-norm at 97.6\%---a 34pp gap using the identical pruning strategy, directly isolating the connector effect.

At $k{=}0.7$ under L2-norm pruning, OpenVLA achieves 51\% while X-VLA achieves 56.2\%. Despite similar absolute values, the degradation \emph{patterns} differ fundamentally: OpenVLA's 51\% represents a smooth continuation from $k{=}0.9$ (63\%, $\Delta = -12$pp), while X-VLA's 56.2\% is a cliff from $k{=}0.9$ (97.6\%, $\Delta = -41$pp).

\subsection{Re-Encoding Mechanism Analysis}

Three lines of evidence confirm that learned re-encoding, rather than compression or a specific architecture, drives VTR's effectiveness on pixel-shuffle:

\begin{enumerate}[leftmargin=*,nosep]
    \item \textbf{M=128 passthrough}: No compression occurs at M=128, yet SR = 58.0$\pm$3.5\% (vs.\ 0\% pruning, 44.8\% baseline). The gain is purely from re-encoding.
    \item \textbf{Architecture control}: MLP adapter (63.4$\pm$2.6\%) $\approx$ cross-attention (60.1$\pm$1.3\%) at M=32, $p = 0.15$. Any learned re-encoding works.
    \item \textbf{Flat M-curve}: M=16 (56.9\%) $\approx$ M=64 (62.0\%) $\approx$ M=128 (58.0\%). Compression ratio is not the explanatory variable; all values cluster within $\pm$5pp.
\end{enumerate}

The M=96 dip (54.5$\pm$3.2\%) is statistically confirmed ($p = 0.034$ vs.\ M=64), likely reflecting a local optimization difficulty at this particular bottleneck dimension rather than a meaningful compression effect.

%% ===== 6. DISCUSSION =====
\section{Discussion}

\textbf{Practical implications.} Our findings provide actionable guidance for VLA deployment. For models with linear projection ($\pi$0-family), aggressive token pruning (up to 90\%) can be applied with no performance cost. For MLP-projector models (OpenVLA-family), practitioners face a smooth quality--efficiency tradeoff: 10\% removal costs $\sim$16\% relative SR, while 50\% removal costs $\sim$49\%. For pixel-shuffle models, pruning is not viable, but VTR provides an alternative pathway to 75\% KV-cache reduction with near-lossless performance. The CV diagnostic enables practitioners to predict which regime applies to a new VLA without expensive closed-loop evaluation.

\textbf{Inference efficiency.} At M=32 (75\% compression), VTR reduces KV-cache memory by 75\% and yields 6\% inference latency speedup, while maintaining 60\%+ SR on SmolVLA.

\textbf{Limitations.} (1) All evaluations use LIBERO simulation; real-robot validation is needed. (2) The CV diagnostic is validated on 4 models (5 including OFT); broader validation across more architectures would strengthen the predictor. (3) RoboFlamingo (Perceiver Resampler connector) is predicted to show intermediate behavior (CV $\sim$0.15) but is not yet validated---this represents a fourth connector family that future work should examine. (4) Statistical power is limited by $n=3$ seeds; no pairwise comparison survives Bonferroni correction at 10 tests. (5) VTR's 22M parameter overhead, while small relative to the base models, adds training cost.

%% ===== 7. CONCLUSION =====
\section{Conclusion}

We present the first systematic analysis of token pruning robustness across VLA architectures, testing 5 models spanning 3 vision connector families. Our findings establish that the effective vision-to-LLM projection pipeline is the primary determinant of token compressibility: pure linear projection is immune (CV = 0.038), linear with deep post-processing shows a threshold-cliff pattern (CV = 0.151), MLP connectors degrade smoothly (CV = 0.215), and pixel-shuffle collapses catastrophically (CV = 0.314). The CV of token importance scores provides a zero-cost, single-forward-pass predictor of this spectrum.

For pixel-shuffle connectors---the hardest case---Vision Token Re-encoding (VTR) recovers from 0\% to 60\%+ SR at 75\% compression, with a flat M-curve confirming that re-encoding itself, not the compression ratio, is the mechanism. VTR's effectiveness is itself connector-dependent (0\% SR on MLP-projector models), reinforcing the centrality of the vision connector in VLA token compressibility.

These findings provide both a diagnostic framework (CV metric) and a targeted solution (VTR for pixel-shuffle) for the emerging challenge of vision token efficiency in VLA deployment.

\bibliographystyle{plain}
\begin{thebibliography}{20}

\bibitem{brohan2023rt2}
A.~Brohan et al.
\newblock RT-2: Vision-Language-Action Models Transfer Web Knowledge to Robotic Control.
\newblock \emph{CoRL}, 2023.

\bibitem{kim2024openvla}
M.~J. Kim et al.
\newblock OpenVLA: An Open-Source Vision-Language-Action Model.
\newblock \emph{arXiv preprint arXiv:2406.09246}, 2024.

\bibitem{black2024pi0}
K.~Black et al.
\newblock $\pi_0$: A Vision-Language-Action Flow Model for General Robot Control.
\newblock \emph{arXiv preprint arXiv:2410.24164}, 2024.

\bibitem{zhen2024xvla}
H.~Zhen et al.
\newblock X-VLA: Cross-Modal Visual Language Action Model.
\newblock \emph{arXiv preprint}, 2024.

\bibitem{smolvla2025}
HuggingFace.
\newblock SmolVLA: A Small Vision-Language-Action Model.
\newblock 2025.

\bibitem{liu2024libero}
B.~Liu et al.
\newblock LIBERO: Benchmarking Knowledge Transfer for Lifelong Robot Learning.
\newblock \emph{NeurIPS}, 2024.

\bibitem{zhai2023sigmoid}
X.~Zhai et al.
\newblock Sigmoid Loss for Language Image Pre-Training.
\newblock \emph{ICCV}, 2023.

\bibitem{lightvla2025}
LightVLA authors.
\newblock LightVLA: Efficient Vision-Language-Action Models via Light Token Pruning.
\newblock \emph{ICRA}, 2026.

\bibitem{seewhatmatters2025}
See What Matters authors.
\newblock See What Matters: Importance-Aware Token Pruning for VLAs.
\newblock \emph{ICML}, 2026.

\end{thebibliography}

\end{document}
